\begin{abstract}
Many real decision problems are partially observed, require simultaneous action
commitment, and exhibit combinatorial structure that causes naive learners to
memorize frequent patterns rather than infer hidden state. This thesis studies
whether causal transformer sequence models can maintain useful beliefs over
hidden opponent attributes from interaction history and whether explicit belief
supervision improves offline policy learning under combinatorial distribution
shift.

We propose a hierarchical tokenization scheme with within-turn self-attention
and cross-turn causal attention, auxiliary belief prediction heads over
active-opponent attributes, and a combinatorial evaluation protocol based on
held-out opponent team keys. Experiments use Metamon Gen~9 OU human replays.
Competitive Pok\'emon is the benchmark domain only. The scientific target is
belief inference and generalization under offline imitation, not ladder win
rate or search-based play strength.

On a 2\,000-battle belief-enriched corpus, turn-only behavioral cloning reaches
20.4\% top-1 action match on random battle holdout and 45.1\% on held-out team
keys. Explicit belief heads reach 28.1\% move-slot accuracy and beat a GRU
baseline, while a latest-turn MLP reaches 35.5\%. Auxiliary belief and matchup
modules do not improve action match or team-level generalization over turn-only
cloning. The opponent-team generalization gap collapses with data scale for
every architecture. These results support nontrivial offline belief learning
and show that belief-to-policy transfer is not automatic at Master's-scale
models and data.

\textbf{Keywords:} partial observability, offline reinforcement learning,
transformers, opponent modeling, combinatorial generalization, behavioral cloning
\end{abstract}

\chapter*{Acknowledgments}
\addcontentsline{toc}{chapter}{Acknowledgments}

I thank my thesis advisor and the research community behind Metamon, poke-env,
and the Pokemon Showdown ecosystem for open tools and datasets that make this
work possible.

\clearpage
